
\documentclass{article}
\usepackage{colortbl}
\usepackage{makecell}
\usepackage{multirow}
\usepackage{supertabular}

\begin{document}

\newcounter{utterance}

\centering \large Interaction Transcript for game `wordle', experiment `medium\_frequency\_words\_no\_clue\_no\_critic', episode 8 with reflexion{-}Qwen3.5{-}27B.
\vspace{24pt}

{ \footnotesize  \setcounter{utterance}{1}
\setlength{\tabcolsep}{0pt}
\begin{supertabular}{c@{$\;$}|p{.15\linewidth}@{}p{.15\linewidth}p{.15\linewidth}p{.15\linewidth}p{.15\linewidth}p{.15\linewidth}}
    \# & \multicolumn{2}{c}{Player} && \multicolumn{2}{c}{Game Master} \\
    \hline

    \theutterance \stepcounter{utterance}  
    & & & \multicolumn{4}{p{0.6\linewidth}}{
        \cellcolor[rgb]{0.9,0.9,0.9}{
            \makecell[{{p{\linewidth}}}]{
                \texttt{\tiny{[P1$\langle$GM]}}
                \texttt{You have attempted this same task before and did not succeed. These are the plans you wrote after those attempts.} \\
\texttt{Trial \#1: New plan: The previous attempt failed because I entered a repetitive loop in the explanation text, repeatedly stating that "hence" was invalid due to the letter 'c' being eliminated, instead of immediately selecting a valid alternative word. This violated the rule to provide a concise explanation and a valid guess. Next time, when I identify that a candidate word contains an eliminated letter, I will immediately discard it and select a new word that strictly adheres to the constraints (e.g., 'r' in position 1, 3, or 5; 'e' in position 4; no 's', 'l', 'a', 't', 'c', 'o', 'y', 'd', 'i'). For instance, I would choose "brine" or "grime" (checking 'g' and 'm' are not eliminated) or "prone" (checking 'p' and 'n') and output the explanation and guess in a single, clean pass without self{-}correction loops.} \\
\\ 
\texttt{Trial \#2: New plan: The previous attempt failed because the explanation text devolved into a stream of consciousness, repeatedly questioning the validity of "rune" and failing to output a final, valid five{-}letter word in the "guess:" field. This violated the rule to provide a concise explanation and a legal guess. Next time, I will strictly separate my internal reasoning from the output. I will first mentally filter the constraints (Pattern: R \_ N \_ \_, 'e' must be in pos 2 or 4, exclude s, l, a, t, i) to identify a valid candidate like "rondo" (if 'o' is allowed) or "riven" (if 'v' is allowed) or "roney" (if 'y' is allowed). Once a valid word is selected, I will write a single, direct sentence explaining the logic (e.g., "Placing 'e' in position 4 to satisfy the yellow feedback while keeping 'r' and 'n' fixed") and immediately output the word in the "guess:" tag without any self{-}correction or hesitation in the text.} \\
\\ 
\\ 
\texttt{You are a language wizard who likes to guess words by using the given rules.} \\
\\ 
\texttt{Welcome to Wordle! You have six attempts to guess the target word, a valid English word of five lowercase letters (a{-}z). Please use the tags "explanation:" and "guess:" to provide a concise explanation for each guess.} \\
\\ 
\texttt{For instance, if your guess is "apple", your response should be} \\
\texttt{explanation: this is a common five{-}letter English word, and I am starting my guess with this word.} \\
\texttt{guess: apple} \\
\\ 
\texttt{After each guess, your answer will be validated, and you will receive feedback indicating which letters are correct (green), which letters are correct but in the wrong position (yellow), and which letters are incorrect (red). This feedback can be useful in determining which letters to include or exclude in your next guess.} \\
\\ 
\texttt{For example, the feedback for "apple" might be:} \\
\texttt{guess\_feedback: a<yellow> p<yellow> p<green> l<yellow> e<red>} \\
\\ 
\texttt{The explanation should contain details about how the guess\_feedback is used to arrive at a new guess.} \\
\\ 
\texttt{Do not generate any other text.} \\
\\ 
\texttt{Let's begin with your first guess.} \\
            }
        }
    }
     \\ \\

    \theutterance \stepcounter{utterance}  
    & \multicolumn{4}{p{0.6\linewidth}}{
        \cellcolor[rgb]{0.561,0.737,0.561}{
            \makecell[{{p{\linewidth}}}]{
                \texttt{\tiny{[P1$\rangle$GM]}}
                \texttt{explanation: I am starting with a common five{-}letter word containing frequent vowels and consonants to gather initial feedback on letter presence and positions.} \\
\texttt{guess: slate} \\
            }
        }
    }
    & & \\ \\

    \theutterance \stepcounter{utterance}  
    & & & \multicolumn{2}{p{0.3\linewidth}}{
        \cellcolor[rgb]{0.9,0.9,0.9}{
            \makecell[{{p{\linewidth}}}]{
                \texttt{\tiny{[GM$|$GM]}}
                \texttt{attempts = 1} \\
\texttt{target\_word = renew} \\
\texttt{guess = slate} \\
\texttt{guess\_feedback = s<red> l<red> a<red> t<red> e<yellow>} \\
            }
        }
    }
    & & \\ \\

    \theutterance \stepcounter{utterance}  
    & & & \multicolumn{4}{p{0.6\linewidth}}{
        \cellcolor[rgb]{0.9,0.9,0.9}{
            \makecell[{{p{\linewidth}}}]{
                \texttt{\tiny{[P1$\langle$GM]}}
                \texttt{guess\_feedback: s<red> l<red> a<red> t<red> e<yellow>} \\
\\ 
\texttt{Provide your response only in this format.} \\
\texttt{explanation: details} \\
\texttt{guess: word} \\
            }
        }
    }
     \\ \\

    \theutterance \stepcounter{utterance}  
    & \multicolumn{4}{p{0.6\linewidth}}{
        \cellcolor[rgb]{0.561,0.737,0.561}{
            \makecell[{{p{\linewidth}}}]{
                \texttt{\tiny{[P1$\rangle$GM]}}
                \texttt{explanation: The feedback indicates that 's', 'l', 'a', and 't' are not in the word, while 'e' is present but not in the fifth position; I will use a new word with common letters and place 'e' in the second position to test its location.} \\
\texttt{guess: crone} \\
            }
        }
    }
    & & \\ \\

    \theutterance \stepcounter{utterance}  
    & & & \multicolumn{2}{p{0.3\linewidth}}{
        \cellcolor[rgb]{0.9,0.9,0.9}{
            \makecell[{{p{\linewidth}}}]{
                \texttt{\tiny{[GM$|$GM]}}
                \texttt{attempts = 2} \\
\texttt{target\_word = renew} \\
\texttt{guess = crone} \\
\texttt{guess\_feedback = c<red> r<yellow> o<red> n<yellow> e<yellow>} \\
            }
        }
    }
    & & \\ \\

    \theutterance \stepcounter{utterance}  
    & & & \multicolumn{4}{p{0.6\linewidth}}{
        \cellcolor[rgb]{0.9,0.9,0.9}{
            \makecell[{{p{\linewidth}}}]{
                \texttt{\tiny{[P1$\langle$GM]}}
                \texttt{guess\_feedback: c<red> r<yellow> o<red> n<yellow> e<yellow>} \\
\\ 
\texttt{Provide your response only in this format.} \\
\texttt{explanation: details} \\
\texttt{guess: word} \\
            }
        }
    }
     \\ \\

    \theutterance \stepcounter{utterance}  
    & \multicolumn{4}{p{0.6\linewidth}}{
        \cellcolor[rgb]{0.561,0.737,0.561}{
            \makecell[{{p{\linewidth}}}]{
                \texttt{\tiny{[P1$\rangle$GM]}}
                \texttt{explanation: The feedback shows 'c' and 'o' are incorrect, while 'r', 'n', and 'e' are in the word but in the wrong positions; since 'e' was yellow in position 5 previously and now yellow in position 5 again, it must be in position 2 or 4, and 'r' and 'n' need to be shifted to valid spots, so I will try "nerve" to place 'e' in position 2 and test 'v'.} \\
\texttt{guess: nerve} \\
            }
        }
    }
    & & \\ \\

    \theutterance \stepcounter{utterance}  
    & & & \multicolumn{2}{p{0.3\linewidth}}{
        \cellcolor[rgb]{0.9,0.9,0.9}{
            \makecell[{{p{\linewidth}}}]{
                \texttt{\tiny{[GM$|$GM]}}
                \texttt{attempts = 3} \\
\texttt{target\_word = renew} \\
\texttt{guess = nerve} \\
\texttt{guess\_feedback = n<yellow> e<green> r<yellow> v<red> e<yellow>} \\
            }
        }
    }
    & & \\ \\

    \theutterance \stepcounter{utterance}  
    & & & \multicolumn{4}{p{0.6\linewidth}}{
        \cellcolor[rgb]{0.9,0.9,0.9}{
            \makecell[{{p{\linewidth}}}]{
                \texttt{\tiny{[P1$\langle$GM]}}
                \texttt{guess\_feedback: n<yellow> e<green> r<yellow> v<red> e<yellow>} \\
\\ 
\texttt{Provide your response only in this format.} \\
\texttt{explanation: details} \\
\texttt{guess: word} \\
            }
        }
    }
     \\ \\

    \theutterance \stepcounter{utterance}  
    & \multicolumn{4}{p{0.6\linewidth}}{
        \cellcolor[rgb]{0.561,0.737,0.561}{
            \makecell[{{p{\linewidth}}}]{
                \texttt{\tiny{[P1$\rangle$GM]}}
                \texttt{explanation: The feedback confirms 'e' is in position 2 (green), while 'n' and 'r' are in the word but not in positions 1 and 3 respectively; 'v' is incorrect. With 'e' fixed at position 2, I need to place 'n' and 'r' in positions 3, 4, or 5, and find a valid word using remaining common letters like 'i' or 'o' for the other slots, so I will try "nerve" again? No, 'v' is red. I will try "nerdy" but 'd' and 'y' are unknown. Let's try "renow" or "reneg". Wait, 'n' was yellow in pos 1, 'r' yellow in pos 3. So 'n' is not 1, 'r' is not 3. 'e' is 2. Possible pattern: \_ e \_ \_ \_. 'n' and 'r' must be in 3, 4, 5. Let's try "renow" (if 'w' is ok) or "reneg". Actually, "nerve" had 'n' at 1 (yellow), 'e' at 2 (green), 'r' at 3 (yellow), 'v' at 4 (red), 'e' at 5 (yellow). So 'n' is not 1, 'r' is not 3. 'e' is 2. 'e' is also in the word (we have two 'e's? Or just one? "nerve" has two 'e's. Feedback: n(y), e(g), r(y), v(r), e(y). This implies the first 'e' is correct, the second 'e' is yellow, meaning there is another 'e' in the word not at pos 5. So the word has two 'e's. One at pos 2, one elsewhere. 'n' and 'r' are in the word. 'v' is out. Let's try "renew". 'r' at 1, 'e' at 2, 'n' at 3, 'e' at 4, 'w' at 5. 'r' was yellow at 3, so 1 is possible. 'n' was yellow at 1, so 3 is possible. 'e' at 4 is possible.} \\
\texttt{guess: renew} \\
            }
        }
    }
    & & \\ \\

    \theutterance \stepcounter{utterance}  
    & & & \multicolumn{2}{p{0.3\linewidth}}{
        \cellcolor[rgb]{0.9,0.9,0.9}{
            \makecell[{{p{\linewidth}}}]{
                \texttt{\tiny{[GM$|$GM]}}
                \texttt{attempts = 4} \\
\texttt{target\_word = renew} \\
\texttt{guess = renew} \\
\texttt{guess\_feedback = r<green> e<green> n<green> e<green> w<green>} \\
            }
        }
    }
    & & \\ \\

    \theutterance \stepcounter{utterance}  
    & & & \multicolumn{2}{p{0.3\linewidth}}{
        \cellcolor[rgb]{0.9,0.9,0.9}{
            \makecell[{{p{\linewidth}}}]{
                \texttt{\tiny{[GM$|$GM]}}
                \texttt{game\_result = WIN} \\
            }
        }
    }
    & & \\ \\

\end{supertabular}
}

\end{document}
